\let\negmedspace\undefined
\let\negthickspace\undefined
\documentclass[journal]{IEEEtran}
\usepackage[a5paper, margin=10mm, onecolumn]{geometry}
%\usepackage{lmodern} % Ensure lmodern is loaded for pdflatex
\usepackage{tfrupee} % Include tfrupee package

\setlength{\headheight}{1cm} % Set the height of the header box
\setlength{\headsep}{0mm}     % Set the distance between the header box and the top of the text

\usepackage{gvv-book}
\usepackage{gvv}
\usepackage{cite}
\usepackage{physics}
\usepackage{amsmath,amssymb,amsfonts,amsthm}
\usepackage{algorithmic}
\usepackage{graphicx}
\usepackage{textcomp}
\usepackage{xcolor}
\usepackage{txfonts}
\usepackage{listings}
\usepackage{enumitem}
\usepackage{mathtools}
\usepackage{gensymb}
\usepackage{comment}
\usepackage[breaklinks=true]{hyperref}
\usepackage{tkz-euclide} 
\usepackage{listings}
% \usepackage{gvv}                                        
\def\inputGnumericTable{}                                 
\usepackage[latin1]{inputenc}                                
\usepackage{color}                                            
\usepackage{array}                                            
\usepackage{longtable}                                       
\usepackage{calc}                                             
\usepackage{multirow}                                         
\usepackage{hhline}                                           
\usepackage{ifthen}                                           
\usepackage{lscape}
\begin{document}

\bibliographystyle{IEEEtran}

\title{9.4.35}
\author{EE25BTECH11022 - sankeerthan}
% \maketitle
% \newpage
% \bigskip
{\let\newpage\relax\maketitle}

\renewcommand{\thefigure}{\theenumi}
\renewcommand{\thetable}{\theenumi}
\setlength{\intextsep}{10pt} % Space between text and floats


\numberwithin{align}{enumi}
\numberwithin{figure}{enumi}
\renewcommand{\thetable}{\theenumi}

\textbf{Question:}A train travels a distance of $480km$ at a uniform speed. If the speed had been 8$\frac{km}{h}$ less, then it would have taken 3 hours more to cover the same distance. We need to find the speed of the train.

\textbf{Solution:}
Let the unknown speed be x.
\begin{align}
\vec{t} =480\myvec{\frac{1}{x-8} \\\frac{1}{x}} \\
\vec{A}^T\vec{X} =3,\text{where}\vec{A} =\myvec{1 \\ -1},
\vec{X} =\myvec{\frac{480}{x-8} \\\frac{480}{x}}
\end{align}
Multiply both sides by $x(x-8)$:
\begin{align}
480x-480(x-8)&=3x(x-8) \\
480x-480x+3840 &=3x^2-24x \\
3x^2-24x-3840 &=0 \\
\myvec{x^2 & x & 1}\myvec{3 \\ -24\\ -3840} =0
\end{align}
Discriminant:
\begin{align}
D = (-24)^2 - 4 \times 3 \times (-3840) = 46656 \\
x = \frac{24 \pm \sqrt{46656}}{6} = \frac{24 \pm 216}{6}
\end{align}
Positive root:
\begin{align}
x = \frac{24 + 216}{6} = 40
\end{align}
Therefore,train speed = $40\frac{km}{h}$

\end{document}
